% GENERATED FILE. Edit canonical lessons, then run npm run generate:books.
% canonical-source-hash: fnv1a64:feebad5c793c2941
% canonical-lessons: TA-C197-ninaivuttu, TA-C197-campati, TA-C197-tavir, TA-C197-kutipeyar, TA-C197-totarpu-kol

\chapter{To remind, To earn, To avoid, To move house, To contact}
\label{ch:ta-to-remind-to-earn-to-avoid-to-move-house-to-contact}
\hlchaptermodality{197}

\begin{chapteropening}
\textbf{By the end of this chapter:} \emph{I can say to remind, to earn, to avoid, to move house and to contact.}

Complete the last lesson of chapter 197: I can say to remind, to earn, to avoid, to move house and to contact.
\end{chapteropening}

\section[\texorpdfstring{\ta{நினைவூட்டு} (niṉaivūṭṭu)}{நினைவூட்டு (niṉaivūṭṭu)}]{\ta{நினைவூட்டு} (niṉaivūṭṭu) — to remind}

\label{lesson:TA-C197-ninaivuttu}



\begin{quote}
Before the new one: say the Tamil for to feel, then the Tamil for to soak.
\end{quote}

\subsection*{You'll want to know: \ta{நினைவூட்டு}}
\textbf{\ta{நினைவூட்டு}} — \emph{niṉaivūṭṭu} — \textquotedblleft{}to remind\textquotedblright{}.

\begin{grammarlens}[title={What you've built}]
One more word for this chapter.
\end{grammarlens}

\subsection*{Your turn}
\begin{itemize}
\raggedright
  \item \emph{Say it:} \emph{niṉaivūṭṭu}
  \item \emph{Say it:} \emph{niṉaivūṭṭu}, once more
  \item \emph{Say it:} say \emph{niṉaivūṭṭu}
  \item \emph{Recall:} say the Tamil for to feel, then the Tamil for to soak, then say \emph{niṉaivūṭṭu} again
\end{itemize}

\begin{tcolorbox}[breakable,colback=teal!4,colframe=teal!35!black,title={Before you move on}]
What does \ta{நினைவூட்டு} mean? (\textquotedblleft{}To remind\textquotedblright{}.) Say it once more.
\end{tcolorbox}
\section[\texorpdfstring{\ta{சம்பாதி} (campāti)}{சம்பாதி (campāti)}]{\ta{சம்பாதி} (campāti) — to earn}

\label{lesson:TA-C197-campati}



\begin{quote}
Before the new one: say the Tamil for to soak, then the Tamil for to remind.
\end{quote}

\subsection*{You'll want to know: \ta{சம்பாதி}}
\textbf{\ta{சம்பாதி}} — \emph{campāti} — \textquotedblleft{}to earn\textquotedblright{}.

\begin{grammarlens}[title={What you've built}]
One more word for this chapter.
\end{grammarlens}

\subsection*{Your turn}
\begin{itemize}
\raggedright
  \item \emph{Say it:} \emph{campāti}
  \item \emph{Say it:} \emph{campāti}, once more
  \item \emph{Say it:} say \emph{campāti}
  \item \emph{Recall:} say the Tamil for to soak, then the Tamil for to remind, then say \emph{campāti} again
\end{itemize}

\begin{tcolorbox}[breakable,colback=teal!4,colframe=teal!35!black,title={Before you move on}]
What does \ta{சம்பாதி} mean? (\textquotedblleft{}To earn\textquotedblright{}.) Say it once more.
\end{tcolorbox}
\section[\texorpdfstring{\ta{தவிர்} (tavir)}{தவிர் (tavir)}]{\ta{தவிர்} (tavir) — to avoid}

\label{lesson:TA-C197-tavir}



\begin{quote}
Before the new one: say the Tamil for to remind, then the Tamil for to earn.
\end{quote}

\subsection*{You'll want to know: \ta{தவிர்}}
\textbf{\ta{தவிர்}} — \emph{tavir} — \textquotedblleft{}to avoid\textquotedblright{}.

\begin{grammarlens}[title={What you've built}]
One more word for this chapter.
\end{grammarlens}

\subsection*{Your turn}
\begin{itemize}
\raggedright
  \item \emph{Say it:} \emph{tavir}
  \item \emph{Say it:} \emph{tavir}, once more
  \item \emph{Say it:} say \emph{tavir}
  \item \emph{Recall:} say the Tamil for to remind, then the Tamil for to earn, then say \emph{tavir} again
\end{itemize}

\begin{tcolorbox}[breakable,colback=teal!4,colframe=teal!35!black,title={Before you move on}]
What does \ta{தவிர்} mean? (\textquotedblleft{}To avoid\textquotedblright{}.) Say it once more.
\end{tcolorbox}
\section[\texorpdfstring{\ta{குடிபெயர்} (kuṭipeyar)}{குடிபெயர் (kuṭipeyar)}]{\ta{குடிபெயர்} (kuṭipeyar) — to move house, to relocate}

\label{lesson:TA-C197-kutipeyar}



\begin{quote}
Before the new one: say the Tamil for to earn, then the Tamil for to avoid.
\end{quote}

\subsection*{You'll want to know: \ta{குடிபெயர்}}
\textbf{\ta{குடிபெயர்}} — \emph{kuṭipeyar} — \textquotedblleft{}to move house, to relocate\textquotedblright{}.

\begin{grammarlens}[title={What you've built}]
One more word for this chapter.
\end{grammarlens}

\subsection*{Your turn}
\begin{itemize}
\raggedright
  \item \emph{Say it:} \emph{kuṭipeyar}
  \item \emph{Say it:} \emph{kuṭipeyar}, once more
  \item \emph{Say it:} say \emph{kuṭipeyar}
  \item \emph{Recall:} say the Tamil for to earn, then the Tamil for to avoid, then say \emph{kuṭipeyar} again
\end{itemize}

\begin{tcolorbox}[breakable,colback=teal!4,colframe=teal!35!black,title={Before you move on}]
What does \ta{குடிபெயர்} mean? (\textquotedblleft{}To move house, to relocate\textquotedblright{}.) Say it once more.
\end{tcolorbox}
\section[\texorpdfstring{\ta{தொடர்பு} \ta{கொள்} (toṭarpu koḷ)}{தொடர்பு கொள் (toṭarpu koḷ)}]{\ta{தொடர்பு} \ta{கொள்} (toṭarpu koḷ) — to contact, to get in touch}

\label{lesson:TA-C197-totarpu-kol}



\begin{quote}
Before the new one: say the Tamil for to avoid, then the Tamil for to move house.
\end{quote}

\subsection*{You'll want to know: \ta{தொடர்பு} \ta{கொள்}}
\textbf{\ta{தொடர்பு} \ta{கொள்}} — \emph{toṭarpu koḷ} — \textquotedblleft{}to contact, to get in touch\textquotedblright{}.

\begin{grammarlens}[title={What you've built}]
One more word for this chapter.
\end{grammarlens}

\subsection*{Your turn}
\begin{itemize}
\raggedright
  \item \emph{Say it:} \emph{toṭarpu koḷ}
  \item \emph{Say it:} \emph{toṭarpu koḷ}, once more
  \item \emph{Say it:} say \emph{toṭarpu koḷ}
  \item \emph{Recall:} say the Tamil for to avoid, then the Tamil for to move house, then say \emph{toṭarpu koḷ} again
\end{itemize}

\begin{tcolorbox}[breakable,colback=teal!4,colframe=teal!35!black,title={Before you move on}]
What does \ta{தொடர்பு} \ta{கொள்} mean? (\textquotedblleft{}To contact, to get in touch\textquotedblright{}.) Say it once more.
\end{tcolorbox}
